\section{R7: library retrieval and type-guided refinement}
\label{sec:retrieval}

\Q{R7.1}{How can dependent and universe-polymorphic types be abstracted usefully?}
\ans Use a staged over-approximation that retains rigid constructors, shared variables, and explicit obligation edges. It should deliberately forget some indices, then refine only the distinctions implicated in failed applications. A graph of result-head names alone is too weak to support the claimed reasoning.

Consider a component of type
\[
 \forall \alpha,\quad \mathrm{List}\ \alpha\to\mathrm{Option}\ \alpha.
\]
The two occurrences of $\alpha$ are the same variable. Replacing them with independent wildcards introduces spurious paths such as a list of naturals becoming an option of Booleans. That may be an acceptable initial over-approximation if failures refine it, but it is wasteful. Preserving repeated-variable relationships is cheap and removes an entire class of spurious paths.

For an indexed component, retain the type constructor and a compact index pattern: a variable, zero, successor, a small affine expression, or unknown. For universes, preserve obvious equalities and inequalities and the distinction between unrestricted elimination and a recursor restricted to \code{Prop}. Do not reject a type merely because the abstraction cannot resolve its universe constraints. It must over-approximate the actual legal instantiations.

Components with multiple premises form \emph{hyperedges}, not ordinary edges: all input obligations must be supplied in a compatible context. Type-class dictionaries and proof prerequisites are additional inputs. A single path through type heads does not account for argument multiplicity, shared substitutions, or linear-style restrictions that Leant2 may impose on a consumed Church eliminator.

A proposed pipeline is: cheap rigid-head retrieval; compatible argument-pattern filtering; relaxed hypergraph reachability; exact dependent application; and a small abstraction refinement after a failed application. The refinement should be reusable, such as a missing index equality or a shared-parameter constraint. Merely rejecting the one concrete path without strengthening the abstraction is not the full type-guided abstraction-refinement loop.

\begin{proposition}[Safe abstract reachability exclusion]
Suppose every concrete well-typed grammar derivation maps to an abstract derivation reaching the abstraction of its result. If the abstract search exhaustively establishes that the requested abstract result is unreachable, then no concrete grammar derivation constructs that result.
\end{proposition}
\begin{proof}
A concrete derivation, if one existed, would map to an abstract derivation by the simulation assumption, contradicting abstract unreachability.
\end{proof}

This is the correct theorem for a retrieval-based exclusion. An incomplete abstract search or an under-approximating index cannot support the premise. In ordinary interactive operation, retrieval should usually order a limited provider pool rather than claim a global impossibility result. Widen the pool fairly or report the pool boundary in the receipt.

A discrimination-tree index is still worthwhile, but ``cost proportional only to query size'' is not a general complexity guarantee. Wildcards, polymorphic providers, and a genuinely large result set can require inspecting many candidates. Measure retrieved-set size, exact-unification attempts, and total retrieval cost under growing libraries, not just the tree lookup itself.

\Q{R7.2}{Has type-guided refinement handled type classes already?}
\ans Yes. Hoogle+ explicitly translates Haskell-style class constraints to dictionary types and dictionary-producing components; instances become ways to construct the required dictionaries. This is described in its implementation section, including examples for equality dictionaries \cite{hoogleplus}. It is a direct published answer to the question, though Lean's dependent and higher-kinded cases require additional work.

Use the same basic representation: a constrained component takes a dictionary as an explicit input in the search model, even if Lean syntax later inserts it implicitly. The difference in Lean is that class parameters may include values, multiple related types, output parameters, universe variables, and local instances. Dictionary resolution can assign metavariables, so it belongs to the shared constraint state.

A cache keyed by only a class name and one type is insufficient. Key it by the fully instantiated class expression, the relevant local-instance environment, transparency and universe dependencies, and assignment version. An unresolved output parameter is part of the interface; a successful dictionary may determine it. The current search already lets instance goals help determine other holes, which should be preserved when retrieval becomes more abstract \cite{searchcore}.

Separate three outcomes: a dictionary found with a checked term, a certified incompatible obligation in a restricted complete class fragment, and ordinary resolution failure or timeout. Only the second is a general exclusion. In most real cases failure means ``not resolved by this procedure under these settings.'' A later local instance, different instantiation, or larger budget can change it.

Do not carry Hoogle+'s evaluation numbers over to Lean's whole library. The relevant question is whether explicit dictionary edges and a refined dependent index reduce exact application attempts at equal recall. A benchmark should include ordinary classes, local instance overrides, multi-parameter classes, and output-parameter cases that force nonlocal assignments.

\Q{R7.3}{Can Lean's selection interface accept data goals and examples?}
\ans At the interface level, data goals are admissible; examples can be sent as optional hints, but there is no guarantee that a registered selector interprets them. The distinction can be established directly from the pinned source rather than guessed from the word ``premise.''

In Lean 4.34.0, \path{src/Lean/LibrarySuggestions/Basic.lean} defines the central type as
\begin{lstlisting}
abbrev Selector := MVarId -> Config -> MetaM (Array Suggestion)
\end{lstlisting}
The configuration has a caller tag, a maximum suggestion count, a filter, and \code{hint : Option String}. The source explicitly says that the hint is arbitrary and may be ignored, and that core Lean provides the interface but no default suggestion engine \cite{leanlibrary}. The interface type does not require the metavariable's target to be a proposition. A specific downstream engine may nevertheless be trained on proof goals, filter out useful definitions, or discard observation features.

A first adapter can use \code{caller := some "leant2.data"}, an appropriate filter for safe component declarations, and a structured textual summary of examples in \code{hint}. This is an integration experiment, not an assertion that an existing model understands the summary. For reliable example-aware retrieval, define a typed Leant2 request with the target, local context, carrier role, allowed components, and observation features, and explicitly implement the translation for each backend.

Maintain separate data-component and proof-lemma retrieval scores. A lemma useful for proving length preservation need not be a good component for computing a list, while a computational definition's body and argument types may matter more than its theorem usage history. The current LeanHammer work is evidence for useful neural selection in dependent proof contexts, not a measured result on carrier or data-component invention \cite{leanhammer}.

For training and evaluation, erase selected definitions from a frozen core, standard-library, and curated user-code environment. Keep only legitimate dependencies available to the synthesizer. Remove the original constant, aliases, generated equation lemmas revealing its implementation, and accidental specification shortcuts. Split by dependency-connected families, not random near-duplicate declarations. Use held-out types and algorithm families as an additional stress test. A corpus with the answer still available under another name primarily measures retrieval leakage.

Retain a no-learning baseline and a provider-oracle track. Retrieval quality should be reported as useful-component recall at $k$, exact application success, and end-to-end certified synthesis rate. A proof-retrieval recall score alone does not answer the user's data-synthesis question.
